\scanpage{232}
\mspage{- 190 -}\hfill IX

\deleted{\unsure{a|s}}sa formation est compatible avec tout changement de traits $S'\longrightarrow S$. Ce dernier fait montre d'ailleurs que, pour définir et étudier (14.1.7), on est ramené par descente finie au cas où $A_K$ a une réduction semi-stable. Dans ce cas, les schémas en groupes $T_K,G_K,B_K,\underline M_K,\ldots$ sont les fibres génériques de schémas en groupes $T,G,B,M\ldots$ sur $S$, $G$ étant une extension du schéma abélien $B$ par le tore $T$ (par. 7). On peut alors considérer le complété formel $\hat G$ de $G$ le long de sa fibre spéciale, et par construction on a un isomorphisme canonique (7.2.6)
\[
 \hat G\xrightarrow{\sim}\handwritten{\widehat}{A^{\unsure{\circ|o}}}\quad, \tag*{(14.1.8)}
\]
d'où par passage aux ``fibres génériques'' (qui sont des groupes ridige-analytiquessur $K$) un isomorphisme canonique
\[
 \hat G^{\mathrm{an}}\xrightarrow{\sim}\handwritten{\widehat}{A^\circ}^{\mathrm{an}}\quad. \tag*{(14.1.9)}
\]
On a d'autre part des homomorphismes canoniques (qui sont des immersions ouvertes rigide-analytiques, mais pas en général des isomorphismes)
\[
 \hat G^{\mathrm{an}}\hookrightarrow G_K^{\mathrm{an}}\quad,\quad
 \handwritten{\widehat}{A^\circ}^{\mathrm{an}}\hookrightarrow A_K^{\mathrm{an}}\quad. \tag*{(14.1.10)}
\]
Ceci dit, l'homomorphisme $e$ de (14.1.7) est caractérisé dé façon unique par la commutativité du diagramme
\diagram{(14.1.11)\qquad
\begin{tikzcd}[ampersand replacement=\&]
 \hat G^{\mathrm{an}}\arrow[r]\arrow[d]\&\handwritten{\widehat}{A^\circ}^{\mathrm{an}}\arrow[d]\\
 G_K^{\mathrm{an}}\arrow[r]\& A_K^{\mathrm{an}}\ ,
\end{tikzcd}}
où les flèches non nommées sont celles de (14.1.9) et (14.1.10).

\handwritten{\unsure{508}}

\mspage{508}
% archive sheet number: 231

